\section{RDF-star Interoperability Demonstration}\label{app:rdf-star}

This appendix demonstrates the concrete interoperability path from ADL Lite EventChains to RDF-star, including export, query, and reversible mapping.

\subsection{EventChain to PROV-O Export}

Consider the EventChain for concept \texttt{disc-capital-trap} with two events: REGISTER (by \texttt{discoverer}) and VALIDATE (by \texttt{reviewer}, confidence 0.85). The PROV-O export is:

\begin{lstlisting}[language=bash]
@prefix prov: <http://www.w3.org/ns/prov#> .
@prefix adl: <https://adl-lite.org/ns/> .

adl:evt-001 a prov:Activity ;
    prov:wasAssociatedWith adl:actor-discoverer ;
    adl:eventType "REGISTER" ;
    adl:previousEvent "genesis" ;
    adl:hash "sha256:abc123..." .

adl:evt-002 a prov:Activity ;
    prov:wasAssociatedWith adl:actor-reviewer ;
    prov:wasInformedBy adl:evt-001 ;
    adl:eventType "VALIDATE" ;
    adl:payloadConfidence 0.85 ;
    adl:hash "sha256:def456..." .
\end{lstlisting}

\subsection{PROV-O to RDF-star Conversion}

RDF-star enables statement-level annotation, which is natural for ADL Lite events, as each event annotates a statement about the concept's lifecycle. The conversion proceeds as follows:

\begin{lstlisting}[language=bash]
# The concept's validation as an RDF-star triple
<< adl:disc-capital-trap adl:status "validated" >>
    adl:validator adl:actor-reviewer ;
    adl:confidence 0.85 ;
    adl:eventHash "sha256:def456..." ;
    adl:timestamp "2024-01-15T14:30:00Z" .

# A relation as RDF-star
<< adl:capital-trap adl:isomorphic-to adl:gradient-explosion >>
    adl:confidence 0.92 ;
    adl:mappingType "topological" ;
    adl:validator adl:actor-reviewer .
\end{lstlisting}

\subsection{SPARQL-star Query}

\begin{lstlisting}[language=bash]
SELECT ?concept ?status ?validator ?confidence
WHERE {
    << ?concept adl:status ?status >>
        adl:validator ?validator ;
        adl:confidence ?confidence .
    FILTER (?confidence > 0.8)
}
\end{lstlisting}

This query retrieves all concepts with validated status and confidence $> 0.8$. It includes the validator identity and confidence value---information that requires reification in standard RDF but is first-class in RDF-star.

\subsection{Reversible Mapping}

The mapping from ADL Lite to RDF-star is \emph{lossy} in one direction: RDF-star discards the cryptographic hash chain (it captures the \emph{semantics} of events but not their \emph{integrity evidence}). A \emph{reversible} mapping is nevertheless possible by extending the RDF-star representation with provenance metadata:

\begin{lstlisting}[language=bash]
<< adl:disc-capital-trap adl:status "validated" >>
    adl:provenance adl:evt-002 .

adl:evt-002 a adl:Event ;
    adl:hash "sha256:def456..." ;
    adl:previousEventHash "sha256:abc123..." .
\end{lstlisting}

With this extension, an ADL Lite system can (1) export to RDF-star for SPARQL querying and Semantic Web interoperability, (2) import from RDF-star by reconstructing the EventChain from \texttt{adl:provenance} links, and (3) verify integrity by checking that the reconstructed chain's hashes match the stored \texttt{adl:hash} values.
